% Netra: concepts explained simply, for viva preparation.
% Build with:  pdflatex concepts.tex   (run twice for the table of contents)
\documentclass[11pt,a4paper]{article}

\usepackage[T1]{fontenc}
\usepackage[utf8]{inputenc}
\usepackage[margin=2.2cm]{geometry}
\usepackage{xcolor}
\usepackage{titlesec}
\usepackage{enumitem}
\usepackage{booktabs}
\usepackage{mdframed}
\usepackage{parskip}
\usepackage[hidelinks]{hyperref}

\definecolor{ink}{HTML}{14213D}
\definecolor{accent}{HTML}{1D6FB8}
\definecolor{panel}{HTML}{EEF3F8}
\definecolor{good}{HTML}{15803D}

\titleformat{\section}{\color{ink}\Large\bfseries}{\thesection}{0.6em}{}
\titleformat{\subsection}{\color{accent}\large\bfseries}{\thesubsection}{0.6em}{}

% A term entry: \term{Name}{plain meaning}{where it shows up in Netra}
\newcommand{\term}[3]{%
  \par\medskip\noindent\textbf{\color{ink}#1}\par
  \nobreak\vspace{1pt}#2\par
  \nobreak\vspace{1pt}\textit{\color{accent}In Netra:} #3\par}

% A box for the one-line answer to give in a viva
\newmdenv[backgroundcolor=panel,linewidth=0pt,skipabove=8pt,skipbelow=8pt,
          innertopmargin=6pt,innerbottommargin=6pt]{saybox}

\newcommand{\qa}[2]{\par\medskip\noindent\textbf{Q. #1}\par\nobreak\vspace{1pt}\textbf{\color{good}A.} #2\par}

\title{\vspace{-1.2cm}\textbf{Netra: Every Concept, Explained Simply}\\[2mm]
\large A plain-language companion for the viva}
\author{Computer Networks Course Case Study}
\date{}

\begin{document}
\maketitle

\begin{saybox}
\textbf{How to use this document.} Each term has three parts: what it means in
plain words, and where it appears in our project. If a teacher points at any part
of the system and asks ``what is this?'', find the term here and read the first
two lines. Section~\ref{sec:viva} has direct answers to the questions most likely
to be asked.
\end{saybox}

\begin{saybox}
\textbf{The project in one sentence.} Netra watches network traffic, decides
whether each connection is normal or an attack using machine learning, blocks the
attacker automatically, and -- this is the part that is ours -- it can also say
``I have never seen this before'' instead of guessing, then group those unknown
connections together and give them a name.
\end{saybox}

\tableofcontents
\newpage

%==============================================================
\section{Computer Networks Concepts}
%==============================================================

\subsection{The basics of how data travels}

\term{Packet}
{The smallest unit of data sent over a network. A big file is cut into many
small packets, each with an address label (header) and a piece of the data
(payload). They travel separately and are reassembled at the other end.}
{Our sniffer captures packets one by one. The signature rules look at individual
packets, which is why they can react in about one second.}

\term{Header and Payload}
{The header is the envelope (who sent it, who receives it, how big it is). The
payload is the letter inside (the actual data).}
{We only read headers and sizes, never the payload. This is why we cannot detect
a SQL injection well: the attack is written inside the payload, and we do not
look there.}

\term{Protocol}
{An agreed set of rules for how two computers talk, so both sides understand each
other. Like a language with grammar.}
{We mainly deal with TCP, UDP and HTTP.}

\term{IP Address}
{The unique number that identifies a computer on a network, for example
\texttt{192.168.10.50}. Like a house address.}
{Every alert names a source IP (the attacker) and a destination IP (the victim).
When we block, we block an IP address.}

\term{Port}
{A numbered door on a computer. One computer runs many services, and each service
listens on its own port: web = 80, secure web = 443, SSH = 22, FTP = 21.}
{A port scan is an attacker knocking on many doors to see which are open. Our
\texttt{port\_scan} rule counts how many different ports one source touches.}

\term{Socket / 5-tuple}
{The five things that uniquely identify one conversation: source IP, source port,
destination IP, destination port, and protocol.}
{We group packets into flows using exactly this 5-tuple.}

\term{TCP (Transmission Control Protocol)}
{A reliable protocol. It checks that everything arrived, in order, and resends
anything lost. It sets up a connection first (the three-way handshake) before
sending data.}
{Most attacks in our dataset are TCP-based.}

\term{UDP (User Datagram Protocol)}
{A fast but unreliable protocol. It just sends, without checking whether anything
arrived. Used for video calls and DNS, where speed matters more than perfection.}
{Present in the traffic, but our main attacks are TCP.}

\term{Three-way handshake}
{How a TCP connection starts: the client sends SYN (``can we talk?''), the server
replies SYN-ACK (``yes, can we?''), the client sends ACK (``yes''). Now the
connection is open.}
{A SYN flood abuses this: the attacker sends thousands of SYNs and never replies
with ACK, so the server keeps thousands of half-open connections waiting.}

\term{TCP flags}
{Single on/off bits in the TCP header that say what a packet is for.
SYN = start, ACK = acknowledge, FIN = finish politely, RST = cut off abruptly,
PSH = deliver this now.}
{Our signature rules read real flags from real packets. Important: we found the
flag counters in the CIC-IDS2017 \emph{dataset} are wrong, so the machine learning
models were \textbf{not} trained on them. This is one of our key findings.}

\term{Flow}
{One complete conversation between two machines: all the packets of the same
5-tuple, grouped together, with statistics such as how long it lasted, how many
packets and bytes went each way, and the gaps between packets.}
{This is the unit our machine learning models score. We turn packets into flows,
then classify each flow.}

\term{Bidirectional flow (forward and backward)}
{A conversation has two directions. Forward is from whoever spoke first;
backward is the reply direction.}
{Many of our features are split this way, for example ``Bwd Packets/s''
(backward packets per second). An attack that sends a lot but receives nothing
back shows up clearly here.}

\term{Inter-Arrival Time (IAT)}
{The gap in time between one packet and the next.}
{A slow attack has very large gaps. ``Flow IAT Mean'' is one of the features that
gives away a slowloris attack.}

\subsection{Network devices and layout}

\term{Switch}
{A device that connects computers inside one local network and forwards frames
only to the correct port.}
{In our Packet Tracer topology, switches connect the campus VLANs.}

\term{Router}
{A device that connects different networks together and chooses the path packets
take between them.}
{Connects our campus, DMZ and WAN segments in the topology.}

\term{Firewall}
{A device or software that allows or blocks traffic based on rules.}
{Our response layer writes real firewall rules using \texttt{iptables}.}

\term{iptables}
{The firewall built into Linux. You give it rules such as ``drop everything from
this IP''.}
{Netra creates its own rule chain called \texttt{NETRA}, so every rule it adds can
be inspected and removed with one command. Rules expire automatically.}

\term{DMZ (Demilitarised Zone)}
{A separate network segment for servers that must be reachable from the internet
(like a web server), kept apart from the internal network so that if the server is
hacked, the attacker is not immediately inside the company.}
{In our topology, the protected web server sits in the DMZ and a Netra sensor
watches it.}

\term{VLAN (Virtual LAN)}
{A way to split one physical switch into several separate logical networks, so
departments cannot see each other's traffic.}
{Used in the campus part of the topology.}

\term{SPAN port / port mirroring}
{A switch feature that copies all traffic passing through it to one monitoring
port, so a security tool can watch without being in the way.}
{This is how an IDS like Netra would be deployed in a real network.}

\term{Promiscuous mode}
{A network card setting that makes it accept all packets it sees, not only those
addressed to it. Required for monitoring.}
{Our sniffer needs this to capture traffic.}

\term{Bandwidth, Throughput and Latency}
{Bandwidth is the maximum capacity of a link. Throughput is how much you actually
get. Latency is the delay before data arrives.}
{A DoS attack attacks all three: it fills the bandwidth, collapses throughput for
real users, and increases latency.}

%==============================================================
\section{Cybersecurity Concepts}
%==============================================================

\subsection{Core ideas}

\term{CIA Triad}
{The three goals of security. \textbf{Confidentiality} (only the right people can
read it), \textbf{Integrity} (nobody changed it secretly),
\textbf{Availability} (the service is up when needed).}
{Our project mainly defends \emph{Availability}, because most attacks we detect
(DoS, DDoS, slowloris) try to take a service down.}

\term{IDS (Intrusion Detection System)}
{A system that watches traffic and raises an alarm when it sees an attack. It only
warns; it does not stop anything.}
{Netra is an IDS plus more.}

\term{IPS (Intrusion Prevention System)}
{Same as an IDS, but it also acts: it blocks the attacker.}
{Netra is an IPS, because it writes firewall rules automatically. This is the ``R''
in NETRA -- Automated \emph{response}.}

\term{Signature-based detection}
{Detect attacks by matching known patterns, like an antivirus matching known
viruses. Very accurate for known attacks, completely blind to new ones.}
{We have four signature rules: SYN flood, port scan, slow DoS, and brute force.
They fire in about one second because they work on packets directly.}

\term{Anomaly-based detection}
{Learn what normal looks like, then flag anything that is different. Can catch new
attacks, but produces more false alarms.}
{Our autoencoder does this. It was trained only on normal traffic and never saw an
attack during training.}

\term{Zero-day attack}
{An attack nobody has seen before, so no signature exists for it and no training
data contains it.}
{This is the central problem our project addresses. We simulate it by removing one
attack type completely from training and then testing against it.}

\term{False Positive and False Negative}
{A false positive is a false alarm: normal traffic reported as an attack. A false
negative is a miss: a real attack that was not reported.}
{False positives are expensive because they block innocent users. Our full model
has a false-positive rate of 0.00112, meaning about 1 in 900 normal flows is
wrongly flagged.}

\term{Defence in depth}
{Do not rely on one protection. Use several layers so that if one fails, another
catches the problem.}
{Netra has four detection layers: rules, two trained classifiers, an autoencoder,
and the open-set test.}

\term{Attack surface}
{Everything an attacker could possibly target: every open port, service and
input.}
{Reducing it (closing unused ports) is a basic defence; a port scan is an attacker
measuring it.}

\subsection{The attacks we detect}

\term{DoS (Denial of Service)}
{One machine overwhelms a server with requests so real users cannot get through.}
{Classes \texttt{DoS} (tools: Hulk, GoldenEye) in our dataset.}

\term{DDoS (Distributed Denial of Service)}
{Same idea, but from many machines at once, so it is far harder to block and much
harder to tell apart from a crowd of real users.}
{Class \texttt{DDoS}. In our test split it is 19,201 flows.}

\term{Slowloris / Slow DoS}
{A clever, quiet DoS. Instead of flooding, it opens many connections and sends
data extremely slowly, keeping every connection alive and using up all the
server's connection slots with almost no bandwidth.}
{Class \texttt{DoS-Slow}. It is the attack our open-set layer detects best,
because its timing pattern is very unusual.}

\term{SYN flood}
{A DoS that sends thousands of SYN packets and never completes the handshake,
filling the server's table of half-open connections.}
{Detected by our \texttt{syn\_flood} signature rule at over 100 SYN per second
from one source.}

\term{Port scan}
{Probing many ports on a target to find which services are running. Usually the
first step of an attack -- reconnaissance.}
{Class \texttt{PortScan}; also caught by a rule at more than 30 distinct ports in
5 seconds.}

\term{Brute force attack}
{Trying username and password combinations repeatedly until one works.}
{Class \texttt{BruteForce} (tools: FTP-Patator and SSH-Patator). Our rule watches
for many HTTP 401 ``Unauthorized'' replies or many new connections to login ports.}

\term{Botnet / Bot}
{A network of infected computers controlled remotely by an attacker. Each infected
machine quietly ``calls home'' to a command-and-control server.}
{Class \texttt{Bot}. It is the hardest class for us, because a bot checking in
looks like ordinary small web traffic.}

\term{Web attack (SQL Injection, XSS, Brute Force)}
{Attacks that abuse a website's input fields. SQL injection tricks the database
into running the attacker's commands; XSS injects scripts into pages other users
view.}
{Class \texttt{WebAttack}. We detect it poorly in live mode and we explain
honestly why: the attack is in the payload text, and we only read flow statistics.}

\term{Infiltration and Heartbleed}
{Two more classes present in CIC-IDS2017. Heartbleed had only 11 examples, far too
few to learn from.}
{We dropped classes with fewer than 100 rows, which is why our final class list has
eight entries.}

\subsection{Response}

\term{Blocking vs Quarantine (rate limiting)}
{Blocking drops all traffic from an address. Quarantine is gentler: it allows a
few packets per second and drops the rest, which slows an attacker without fully
cutting off a possibly innocent user.}
{High severity blocks for 120 seconds; medium severity quarantines at 10 packets
per second for 60 seconds; low severity is only recorded.}

\term{Timed and reversible response}
{Every automatic block expires by itself after a set time and can be undone
manually.}
{This is our main safety argument: the worst case of a false positive is a short,
bounded outage for one address, not a permanent ban.}

\term{Never-block list (allow list)}
{A list of addresses that must never be blocked, no matter what.}
{Protects loopback, the gateways and the protected server itself, so the system
cannot cut off its own network.}

%==============================================================
\section{Machine Learning Concepts}
%==============================================================

\subsection{The fundamentals}

\term{Machine Learning}
{Instead of writing rules by hand, you show a program many labelled examples and it
works out the rules itself.}
{We show it 635,280 example flows labelled Benign, DoS, DDoS and so on, and it
learns to tell them apart.}

\term{Feature}
{One measurable property of a thing, used as input to the model. For a flow:
duration, packets per second, average packet size, and so on.}
{We use 59 features in the full model and 6 in the lite model.}

\term{Label / Class}
{The correct answer for a training example.}
{Our eight classes: Benign, Bot, BruteForce, DDoS, DoS, DoS-Slow, PortScan,
WebAttack.}

\term{Supervised learning}
{Learning from examples that come with the correct answer.}
{Our Random Forest and XGBoost are supervised.}

\term{Unsupervised learning}
{Learning from data with no answers given; the model finds structure by itself.}
{Our autoencoder and our family clustering are unsupervised.}

\term{Training, Validation and Test sets}
{Split the data in three. \textbf{Train} teaches the model. \textbf{Validation}
tunes choices and picks the best model. \textbf{Test} is touched only at the very
end, to get an honest score.}
{Train 635,280 flows, validation 329,392, test 329,392. We never trained on the
test set.}

\term{Overfitting}
{When a model memorises the training data instead of learning the general pattern.
It scores brilliantly on data it has seen and badly on anything new.}
{We guard against it with a separate test set -- and our cross-dataset experiment
shows a deeper version of this problem, where the model learned the \emph{network}
rather than the attack.}

\term{Feature scaling / Normalisation}
{Putting all features on a comparable range, so a feature measured in millions
does not drown out one measured in single digits.}
{We apply \texttt{sign(x)*log(1+|x|)} first, to squash very large values, then
standardise. The scaler is fitted on the training data only.}

\term{Class imbalance}
{When one class has far more examples than the others. Benign traffic hugely
outnumbers attacks, so a lazy model could call everything Benign and still be
99\% ``accurate''.}
{We handle it with balanced class weights, and we report \emph{macro} F1, which
gives every class equal importance regardless of size.}

\subsection{The models we used}

\term{Decision Tree}
{A flowchart of yes/no questions. ``Is packets per second above 500? If yes, ask
the next question; if no, call it Benign.''}
{The building block of our main model.}

\term{Random Forest}
{Many decision trees, each trained on a random part of the data and random
features. They vote, and the majority wins. Much more reliable than one tree.}
{Our selected model, with 100 trees. It won on validation macro F1 for both
feature sets.}

\term{XGBoost (Gradient Boosting)}
{Also many trees, but built one after another, where each new tree tries to fix
the mistakes of the ones before it.}
{We trained it as a competitor. It scored very slightly lower on macro F1, so the
Random Forest was selected.}

\term{Autoencoder}
{A neural network shaped like an hourglass. It is trained to squeeze data into a
small summary and then rebuild the original from it. Trained only on normal
traffic, it rebuilds normal traffic well and unusual traffic badly.}
{Our shape is 59-32-16-8-16-32-59. The rebuilding error is the anomaly score: a
high error means ``this does not look like anything normal''.}

\term{Reconstruction error}
{How different the autoencoder's rebuilt version is from the original input.}
{If it goes above a threshold (the 99th percentile of errors on normal traffic),
we call the flow anomalous.}

\term{Threshold}
{The cut-off that turns a score into a yes/no decision.}
{Our autoencoder threshold is 0.1463, taken from the dataset. We also wrote a
\texttt{calibrate} command that re-derives it from the real network being watched.}

\subsection{Measuring how good a model is}

\term{Confusion Matrix}
{A table showing, for each true class, what the model actually predicted. The
diagonal is correct; everything off the diagonal is a mistake.}
{We generate one for each model in \texttt{reports/}.}

\term{Accuracy}
{The share of all predictions that were correct. Misleading when classes are
imbalanced.}
{Ours is 0.9987, but we do not lead with it, for exactly that reason.}

\term{Precision}
{Of everything the model called an attack, how much really was an attack.
``When it shouts, is it right?''}
{Low precision means many false alarms.}

\term{Recall (Detection Rate)}
{Of all the real attacks, how many the model caught. ``Does it miss things?''}
{Our full model's detection rate is 0.9980.}

\term{F1 score}
{A single number balancing precision and recall (their harmonic mean). Useful when
you care about both.}
{Our headline number is \textbf{macro} F1 = 0.9570.}

\term{Macro F1 vs Weighted F1}
{Macro averages the F1 of each class equally, so a tiny class counts as much as a
huge one. Weighted averages by class size, so big classes dominate.}
{We select models on macro F1 so that rare attacks are not ignored.}

\term{ROC-AUC}
{A score from 0.5 (useless, same as guessing) to 1.0 (perfect) measuring how well
a model separates two groups across every possible threshold.}
{Our autoencoder reaches 0.9603.}

\term{False Positive Rate (FPR)}
{The share of normal traffic wrongly flagged.}
{0.00112 for the full model. On 278,975 benign test flows this is 313 false
alarms.}

%==============================================================
\section{The New Ideas in Our Project}
%==============================================================

\begin{saybox}
This section is the part that makes our project different from a standard
``train a classifier on CIC-IDS2017'' project. Expect questions here.
\end{saybox}

\term{Closed-set classification (the problem)}
{A trained classifier can only answer with classes it was taught. Ask it about
something new and it must still pick one of its eight answers -- it has no option
for ``none of these''. It is like a multiple-choice exam with no ``none of the
above''.}
{We proved the damage: when we remove an attack class from training and then show
it to the model, in 6 out of 7 cases the model calls most of that attack
\emph{Benign} and says nothing at all.}

\term{Open-set recognition (our solution)}
{Give the system a ninth possible answer: ``this matches nothing I was trained
on''.}
{Implemented in \texttt{src/models/openset.py}. It sits after the classifiers:
if they say Benign, we double-check whether the flow actually resembles anything
we have seen.}

\term{Nearest-neighbour distance}
{To judge whether a new flow resembles class X, measure how far it is from the
10 closest training examples of class X. Close means familiar; far means strange.}
{This is our ``non-conformity score''.}

\term{Conformal p-value}
{A fair way to turn a distance into a probability-like number between 0 and 1. It
asks: ``among genuine members of this class, what fraction were at least this
strange?'' A p-value of 0.5 means perfectly typical; 0.001 means extremely
unusual.}
{The key benefit: if we reject everything below a level $\alpha$, then exactly
$\alpha$ of genuinely normal traffic gets rejected. The cost is known \emph{in
advance} instead of being tuned by trial and error.}

\term{Alpha ($\alpha$), the significance level}
{The dial that decides how suspicious we need to be. $\alpha = 0.05$ means we
accept sending 5\% of known traffic to the ``unknown'' path.}
{We promised 0.05 and measured 0.0468 on the test data. The promise held at every
level we tried, from 0.01 to 0.5.}

\term{Mahalanobis distance and whitening}
{Ordinary straight-line distance is unfair when some features vary much more than
others. Mahalanobis distance accounts for how spread out the data is in each
direction. ``Whitening'' transforms the space so that ordinary distance becomes
this fairer distance.}
{We use it so that all 59 features count sensibly.}

\term{Why a single Gaussian failed (an honest negative result)}
{The textbook method models each class as one bell-shaped cloud. Normal network
traffic is not one cloud -- it is several separate clouds (web browsing, DNS, file
transfer, idle). One big ellipse drawn around them also covers the empty space
between them, so strange traffic lands in that empty space and is accepted.}
{Measured: on the UNSW-NB15 dataset, this method rejected only 0.22\% of a
completely foreign network. Our nearest-neighbour method rejected 95.09\%.}

\term{Clustering}
{Grouping similar things together without being told the groups in advance.}
{We cluster the rejected flows so that one new attack becomes one object the
operator can look at, instead of hundreds of separate alarms.}

\term{Cosine distance (direction, not size)}
{Measures whether two things point the same way, ignoring how big they are.}
{We group flows by the \emph{direction} in which they differ from normal. This
means the same attack run gently and run aggressively lands in the same family.
Using plain distance instead broke one attack into 24 separate groups.}

\term{Unknown family}
{A cluster of unrecognised flows, given a permanent name like
\texttt{NETRA-U001}, with a fingerprint of what makes it odd and the closest class
we do know.}
{An alert now reads: \emph{``NETRA-U001: 47 unrecognised flows from 3 sources.
Flow IAT Mean = 10.5 s, no backward packets. Nearest known class DoS-Slow,
p = 0.004''} -- instead of the old, useless \emph{``Anomaly (unknown)''}.}

\term{Promotion (the human in the loop)}
{An analyst looks at a family, decides what it is, and gives it a real name. Every
future alert about that family uses that name, and the name survives restarts.}
{This is deliberate: the computer can prove something is new and describe it, but
only a human can say what it is.}

\term{Leave-one-class-out testing}
{To test fairly whether we detect unknown attacks, remove one attack class
completely -- from training, from the scaling, from the calibration -- then show it
to the system. It has genuinely never seen it.}
{We repeat this for all seven attack classes. This is our zero-day experiment.}

\term{Cross-dataset evaluation}
{Train on one dataset, test on a completely different one captured by other
researchers on another network.}
{Our most important negative result: detection falls from 0.9976 to 0.0003,
because the model learned the packet rates of \emph{that particular network},
not what an attack looks like.}

\term{Domain shift}
{When the data a model meets in the real world is different from what it was
trained on, so its answers quietly become wrong.}
{Our open-set layer detects this: on UNSW-NB15 it reports 95.09\% of traffic as
unrecognised, versus 4.68\% at home. That is a sensor correctly saying ``I do not
belong here'', which a normal classifier can never say.}

\term{Explainability (Saabas decomposition)}
{Rather than just saying ``attack'', the system shows which features caused that
decision. For a Random Forest this can be computed exactly, by following the path
each tree took.}
{Our alerts read: \emph{``DoS-Slow predicted with confidence 0.97, driven by
Bwd Packets/s = 0.0, Flow Packets/s = 0.111, Flow IAT Mean = 10500000''}.}

%==============================================================
\section{The Numbers You Should Know}
%==============================================================

\begin{table}[h]
\centering
\small
\begin{tabular}{@{}lll@{}}
\toprule
\textbf{Thing} & \textbf{Number} & \textbf{What it means} \\
\midrule
Dataset & CIC-IDS2017 & 2,830,743 flows loaded; 2,195,946 after cleaning \\
Classes & 8 & Benign + 7 attack types \\
Features & 59 full, 6 lite & lite = computable live from packets \\
Split & 635,280 / 329,392 / 329,392 & train / validation / test \\
Full model & Random Forest & macro F1 \textbf{0.9570}, accuracy 0.9987 \\
Detection rate & 0.9980 & 50,314 of 50,417 attack flows caught \\
False-positive rate & 0.00112 & 313 of 278,975 benign flows \\
Lite model & macro F1 0.8400 & the one that runs on live packets \\
Autoencoder & ROC-AUC 0.9603 & never saw an attack in training \\
Cross-dataset & 0.9976 $\rightarrow$ 0.0003 & the model does not transfer \\
Open-set promise & $\alpha$ = 0.05 & measured 0.0468 on known traffic \\
Zero-day (DoS-Slow) & 0.835 $\rightarrow$ 0.016 & share of attack nobody would see \\
Domain shift & 95.09\% vs 4.68\% & unknown rate abroad vs at home \\
Families & 5 of 7 found & purity 0.991, 16 false families \\
Tests & 96 passing & automated test suite \\
\bottomrule
\end{tabular}
\end{table}

%==============================================================
\section{Questions a Teacher May Ask}
\label{sec:viva}
%==============================================================

\qa{What exactly does your project do?}
{It watches network traffic, converts packets into flows, and scores each flow
with signature rules, two machine learning classifiers and an autoencoder. If it
finds an attack it automatically writes a temporary firewall rule to block the
source. On top of that, it can report that a flow matches none of the attack types
it knows, group such flows into named families, and let an analyst name them.}

\qa{Why machine learning? The rules already catch the attacks in your demo.}
{Rules only catch attacks someone already thought of and wrote down. The models
classify flows the rules say nothing about, and the autoencoder needs no attack
labels at all. The demo deliberately uses attacks the rules cover, because that
makes it fire in one second in front of an audience.}

\qa{What is the difference between IDS and IPS? Which is yours?}
{An IDS only detects and alerts. An IPS also acts to stop the attack. Ours is an
IPS, because it writes real iptables rules. Every rule is timed and reversible.}

\qa{Why did you use a Random Forest instead of deep learning?}
{We compared it against XGBoost on validation macro F1 and Random Forest won. It
also trains in about 31 seconds, scores a flow in 2.5 microseconds, and -- most
usefully -- its decisions can be explained exactly, feature by feature. A deep
network would be slower and harder to justify to an operator.}

\qa{Your accuracy is 99.87\%. Is the problem solved?}
{No, and that number is the least interesting one in the report. It describes
CIC-IDS2017, not intrusion detection. We tested the same model on UNSW-NB15 and
detection collapsed from 99.76\% to 0.03\%, because it had learned the packet
rates of one particular network. Reporting only in-domain accuracy hides this.}

\qa{What is the novelty? Isn't this just a classifier on a public dataset?}
{The novelty is that our system can refuse to answer. A classifier always returns
one of its trained classes, so a new attack is silently labelled Benign -- we
measured this happening in 6 of 7 cases. We added a conformal open-set test whose
false-alarm cost is fixed mathematically in advance, and an online clusterer that
turns the rejected flows into named families with readable fingerprints.}

\qa{How can you set a threshold without having seen the new attacks?}
{That is exactly why we used a conformal test. The threshold is a significance
level: $\alpha$ is the fraction of genuinely normal traffic that will be rejected,
guaranteed by construction. So we choose it from our false-alarm budget, not by
fitting to attack data. We promised 0.05 and measured 0.0468.}

\qa{How do you know a discovered family is real and not random noise?}
{Three ways. Purity across the folds is 0.991, meaning families do not mix the new
attack with normal traffic. The largest family's precision averages 0.909. And the
families turned out to match the individual attack \emph{tools}: one family was
274 FTP-Patator flows, another was 2,639 Hulk flows with no GoldenEye at all --
a distinction finer than the labels we trained on.}

\qa{What stops it blocking innocent users?}
{Four things: weak evidence needs three flagged flows from the same source within
10 seconds; an unknown family must grow to 50 flows before it is reported; every
block expires on its own; and a never-block list protects critical addresses.}

\qa{What is your biggest weakness?}
{We never read packet contents, so web attacks are nearly invisible to us. We
measured this precisely: held-out WebAttack flows are ranked as \emph{more} normal
than normal traffic, because a SQL injection at the flow-statistics level is just
a small HTTP request. No threshold can fix that; only payload inspection can.}

\qa{Did anything fail during the project?}
{Several things, and they are in the report. The dataset's TCP flag counters are
broken, so we removed all twelve. The model does not transfer to another network.
A single-Gaussian open-set test rejected nothing. Plain-distance clustering broke
one attack into 24 pieces. And our first idea for deciding which families to
report -- demanding a stronger statistical rejection -- reported only 1 of 7 real
attacks, so we replaced it with a size threshold chosen from a measured curve.}

\qa{What would you do next?}
{Add payload inspection for web attacks, feed named families back into training so
each discovery becomes a new class, and correlate alerts into attack campaigns so
a multi-stage intrusion reads as one story instead of dozens of alerts.}

\vfill
\begin{saybox}
\textbf{If you only remember one thing:} every intrusion detection system can tell
you what it recognises. Ours can also tell you, in a mathematically controlled
way, when it recognises \emph{nothing} -- and then describe what it found so a
person can name it.
\end{saybox}

\end{document}
